\documentclass[12pt,a4paper]{article}

\usepackage[utf8]{inputenc}
\usepackage[vietnamese]{babel}
\usepackage{amsmath,amssymb,amsthm}
\usepackage{geometry}
\usepackage{enumitem}
\usepackage[most]{tcolorbox}
\usepackage{hyperref}
\usepackage{fancyhdr}
\usepackage{tikz}
\usepackage{eso-pic}
\usepackage{xcolor}

\geometry{a4paper, margin=2cm, bottom=2.5cm}
\hypersetup{colorlinks=true, linkcolor=blue!70!black}
\setlist[itemize]{topsep=4pt, itemsep=2pt}
\setlist[enumerate]{topsep=4pt, itemsep=2pt}

\AddToShipoutPictureFG{%
  \AtPageCenter{%
    \begin{tikzpicture}[remember picture, overlay]
      \node[rotate=45, scale=1, opacity=0.06]{\fontsize{60}{72}\selectfont\bfseries\sffamily Đặng Tấn Quí};
    \end{tikzpicture}%
  }%
}

\pagestyle{fancy}
\fancyhf{}
\fancyhead[L]{\small\textit{Đại số tuyến tính nâng cao}}
\fancyhead[R]{\small\textit{Đặng Tấn Quí}}
\fancyfoot[C]{\thepage}
\fancyfoot[L]{\footnotesize\textcopyright\ Đặng Tấn Quí}
\fancyfoot[R]{\footnotesize SĐT: 0377\,122\,210}
\renewcommand{\headrulewidth}{0.4pt}
\renewcommand{\footrulewidth}{0.4pt}

\fancypagestyle{plain}{\fancyhf{}\fancyfoot[C]{\thepage}\fancyfoot[L]{\footnotesize\textcopyright\ Đặng Tấn Quí}\fancyfoot[R]{\footnotesize SĐT: 0377\,122\,210}\renewcommand{\headrulewidth}{0pt}\renewcommand{\footrulewidth}{0.4pt}}

\newtcolorbox{dinhnghia}[1][]{enhanced, breakable, colback=blue!3!white, colframe=blue!60!black, fonttitle=\bfseries, title={Định nghĩa}, #1}
\newtcolorbox{dinhly}[1][]{enhanced, breakable, colback=green!3!white, colframe=green!50!black, fonttitle=\bfseries, title={Định lý}, #1}
\newtcolorbox{tinhchat}[1][]{enhanced, breakable, colback=yellow!5!white, colframe=orange!50!black, fonttitle=\bfseries, title={Tính chất}, #1}
\newtcolorbox{phuongphap}[1][]{enhanced, breakable, colback=gray!5!white, colframe=gray!50!black, fonttitle=\bfseries, title={Phương pháp}, #1}
\newtcolorbox{luuy}[1][]{enhanced, breakable, colback=red!3!white, colframe=red!50!black, fonttitle=\bfseries, title={Lưu ý}, #1}
\newtcolorbox{congthuc}[1][]{enhanced, breakable, colback=cyan!3!white, colframe=cyan!50!black, fonttitle=\bfseries, title={Công thức}, #1}
\newtcolorbox{vidu}[1][]{enhanced, breakable, colback=violet!3!white, colframe=violet!50!black, fonttitle=\bfseries, title={Ví dụ}, #1}

\begin{document}

\thispagestyle{empty}
\begin{tikzpicture}[remember picture, overlay]
  \fill[blue!80!black] (current page.south west) rectangle (current page.north east);
  \node[anchor=west, white, font=\fontsize{22}{28}\bfseries\selectfont]
    at ([xshift=3cm, yshift=-7.5cm]current page.north west) {ĐẠI SỐ TUYẾN TÍNH NÂNG CAO};
  \node[anchor=west, white, font=\fontsize{18}{24}\selectfont]
    at ([xshift=3cm, yshift=-9cm]current page.north west) {Không gian vectơ, ánh xạ tuyến tính --- Năm 1, HK2};
  \node[anchor=west, white, font=\Large\bfseries] at ([xshift=3cm, yshift=-13cm]current page.north west) {Biên soạn:};
  \node[anchor=west, yellow!80!white, font=\fontsize{22}{28}\bfseries\selectfont] at ([xshift=3cm, yshift=-14.5cm]current page.north west) {ĐẶNG TẤN QUÍ};
  \node[anchor=south east, white, font=\Large\bfseries] at ([xshift=-2cm, yshift=3cm]current page.south east) {\the\year};
\end{tikzpicture}

\newpage
\tableofcontents
\newpage

\section{Không gian vectơ}

\subsection{Định nghĩa và ví dụ}

\begin{dinhnghia}
  \textbf{Không gian vectơ} $V$ trên trường $\mathbb{K}$: tập $V$ với phép cộng $+$ và nhân vô hướng $\cdot$ thỏa 8 tiên đề (đóng, giao hoán, kết hợp, phần tử không, phần tử đối, phân phối, đơn vị).
\end{dinhnghia}

\begin{vidu}
  $\mathbb{R}^n$, $M_{m \times n}(\mathbb{R})$, $\mathbb{R}[x]$ (đa thức), $C[a,b]$ (hàm liên tục).
\end{vidu}

\subsection{Không gian con}

\begin{dinhnghia}
  $W \subset V$ là \textbf{không gian con} nếu $W$ đóng với phép cộng và nhân vô hướng.
\end{dinhnghia}

\begin{phuongphap}[title={Kiểm tra không gian con}]
  $W \ne \varnothing$ và $\forall u, v \in W$, $\forall \lambda \in \mathbb{K}$: $u + v \in W$, $\lambda u \in W$.
\end{phuongphap}

\subsection{Tổ hợp tuyến tính, sinh, cơ sở}

\begin{dinhnghia}
  $v$ là \textbf{tổ hợp tuyến tính} của $\{v_1, \ldots, v_k\}$ nếu $v = \lambda_1 v_1 + \cdots + \lambda_k v_k$.
\end{dinhnghia}

\begin{dinhnghia}
  $\mathrm{span}\{v_1, \ldots, v_k\}$: tập mọi tổ hợp tuyến tính. Hệ \textbf{sinh} $V$ nếu $\mathrm{span}\{v_1, \ldots, v_k\} = V$.
\end{dinhnghia}

\begin{dinhnghia}
  \textbf{Cơ sở} của $V$: hệ vừa độc lập tuyến tính vừa sinh $V$. Mọi cơ sở có cùng số phần tử $= \dim V$.
\end{dinhnghia}

\section{Ánh xạ tuyến tính}

\begin{dinhnghia}
  $T: V \to W$ \textbf{tuyến tính} nếu $T(\alpha u + \beta v) = \alpha T(u) + \beta T(v)$.
\end{dinhnghia}

\begin{dinhnghia}
  \textbf{Núcleus (hạt nhân)}: $\ker T = \{v \in V \mid T(v) = 0\}$. \textbf{Ảnh}: $\mathrm{Im}\,T = T(V)$.
\end{dinhnghia}

\begin{dinhly}[title={Định lý chiều}]
  $\dim V = \dim(\ker T) + \dim(\mathrm{Im}\,T)$.
\end{dinhly}

\begin{dinhnghia}
  \textbf{Ma trận của $T$} trong cơ sở $\mathcal{B}_V$, $\mathcal{B}_W$: $[T(v)]_{\mathcal{B}_W} = A [v]_{\mathcal{B}_V}$.
\end{dinhnghia}

\section{Trị riêng và chéo hóa}

\begin{dinhnghia}
  $\lambda$ là \textbf{trị riêng} của $A$ nếu $\exists x \ne 0$: $Ax = \lambda x$. $x$: vectơ riêng.
\end{dinhnghia}

\begin{congthuc}
  $\lambda$ trị riêng $\Leftrightarrow$ $\det(A - \lambda I) = 0$ (phương trình đặc trưng).
\end{congthuc}

\begin{dinhnghia}
  $A$ \textbf{chéo hóa được} nếu tồn tại $P$ khả nghịch: $P^{-1}AP = D$ (ma trận chéo). Khi đó các cột của $P$ là vectơ riêng.
\end{dinhnghia}

\begin{dinhly}
  $A$ chéo hóa được $\Leftrightarrow$ $A$ có $n$ vectơ riêng độc lập tuyến tính $\Leftrightarrow$ tổng số chiều các không gian riêng $= n$.
\end{dinhly}

\section{Dạng toàn phương}

\begin{dinhnghia}
  \textbf{Dạng toàn phương} $Q(x) = x^T A x = \sum_{i,j} a_{ij} x_i x_j$, $A$ đối xứng.
\end{dinhnghia}

\begin{dinhly}[title={Luật quán tính Sylvester}]
  Số hệ số dương và âm trong dạng chính tắc không phụ thuộc cách đưa về chính tắc.
\end{dinhly}

\begin{dinhnghia}
  $Q$ \textbf{xác định dương} nếu $Q(x) > 0$ với mọi $x \ne 0$. Tương tự xác định âm, nửa xác định.
\end{dinhnghia}

\end{document}
